% 316_Riemann_Zeta_De_ch.tex -- Kapitelinhalt Dok. 316
% FFGFT und die Riemannsche Zeta-Funktion: Reichweite und Grenze

\section*{Vorbemerkung: die Antwort in Kurzform}
\addcontentsline{toc}{section}{Vorbemerkung: die Antwort in Kurzform}

\begin{keyresult}[Ergebnis]
Die Riemann-Hypothese ist eine \textbf{arithmetische} Angelegenheit.
Der Torus und $\xi$ tragen zu ihrer L\"osung nichts bei. Das Ergebnis
dieses Dokuments ist ein \emph{negatives} -- drei naheliegende
geometrische Zug\"ange werden mit Begr\"undung ausgeschlossen.
\end{keyresult}

Die pauschale Fassung w\"are an einer Stelle zu grob, daher die
Pr\"azisierung.

\textbf{Was reine Zahlensache bleibt.} Wo die Nullstellen liegen, wie
sie zueinander stehen, wie man sie einzeln fasst -- das tr\"agt die
Arithmetik, die Primzahlen. Drei geometrische Zug\"ange wurden gepr\"uft
und sind alle negativ: keine kompakte Geometrie kann die Nullstellen als
Spektrum tragen (Abschnitt~\ref{sec:weyl}, Satzebene), $\xi$ verschiebt
die kritische Linie nicht (Abschnitt~\ref{sec:casimir}), und die
$\xi$-Leiter diskretisiert nichts (Abschnitt~\ref{sec:rand}, Test
out-of-sample negativ).

\textbf{Was dennoch nicht blo\ss{} Zahlensache ist.} Zweierlei. Erstens
ist die Riemann-Zeta ein \emph{Faktor} im Spektrum des
$\mathbb{Z}^4$-Torus -- eine Identit\"at, kein Bild
(Abschnitt~\ref{sec:faktor}). Die Riemann-Hypothese ist damit eine
Aussage \"uber ein geometrisches Objekt des Rahmens; nur eben eine, die
die Geometrie nicht beantworten kann. Zweitens sagt die Form der
Nullstellendichte, \emph{wohin} die Frage geh\"ort: $N(T)$ w\"achst wie
$T\ln T$, und das kann kein Laplace auf irgendeiner kompakten
Mannigfaltigkeit liefern -- es verlangt eine Dilatationsstruktur, also
die ausgerollte Dom\"ane. Die Geometrie zeigt die Adresse an, ohne die
Antwort zu liefern.

\textbf{Das Merkw\"urdige ist gerade, dass beides zusammen auftritt:}
Die Zahlen sitzen nachweislich im Torusspektrum, aber der Torus kann sie
nicht erzeugen.

\textbf{Der Ertrag.} Drei Sackgassen sind geschlossen, mit Begr\"undung
statt mit Vermutung. Zus\"atzlich ist eine eigene Fehlkonstruktion
dokumentiert: ein zun\"achst tragf\"ahig wirkendes Zellargument erwies
sich beim Ausformulieren als Umschreibung von $\hbar = 1$
(Abschnitt~\ref{sec:zelle}) und wurde zur\"uckgenommen.

\section{Ausgangslage und Anspruch}

\textbf{[Q]} Der Korpus erw\"ahnt die Riemann-Hypothese an genau einer
Stelle: Dok.~024 notiert im KI-/Netzwerk-Kontext, der
\enquote{Zahlenraum} habe Dimensionen, die \enquote{durch die
Riemann-Hypothese und $\zeta$-Funktionen moduliert werden}. Das ist eine
konzeptuelle Randbemerkung ohne epistemischen Marker und ohne
Herleitung. Eine ausgearbeitete Verbindung existiert im Korpus nicht.

\textbf{[K]} Zugleich enth\"alt Dok.~314 Material, das die
$\zeta$-Funktion unmittelbar ber\"uhrt: Epstein-Zeta, Casimir-Punkt,
W\"armekern auf $\text{T}^4$. Dieses Dokument pr\"uft systematisch, was
daraus folgt.

\textbf{Anspruch.} Es beweist nichts \"uber die Riemann-Hypothese und
behauptet das auch nicht. Gekl\"art werden drei Dinge: welche
Verbindungen als Identit\"aten bestehen \textbf{[K]}, welcher Weg
konstruktiv tr\"agt \textbf{[B]}, und welche naheliegenden Wege
\emph{beweisbar} nicht gangbar sind \textbf{[X]}. Der Ausschluss ist
dabei der wertvollste Teil -- er erspart Arbeit an Sackgassen.

\section{Zwei Befunde, die als Identit\"aten bestehen}
\label{sec:faktor}

\subsection{Die Riemann-Zeta ist ein Faktor der Torus-Spektral-Zeta}

\textbf{[K]} Dok.~314 f\"uhrt die geschlossene Form der
$\mathbb{Z}^4$-Spektral-Zeta (dort dreifach verifiziert; hier gegen die
Direktsumme mit analytischer Schwanzkorrektur auf $10^{-7}$ bis
$10^{-12}$ best\"atigt):
\[
  Z_{\mathbb{Z}^4}(s) \;=\; 8\,\bigl(1 - 4^{1-s}\bigr)\,
  \zeta(s)\,\zeta(s-1).
\]
Das ist keine Analogie, sondern eine Identit\"at. Die Nullstellen der
Torus-Spektral-Zeta haben damit drei Quellen:

\begin{center}
\begin{tabular}{ll}
\toprule
Quelle & Ort der Nullstellen \\
\midrule
$\zeta(s) = 0$ & $\Re(s) = 1/2$ (nichttrivial -- die RH) \\
$\zeta(s-1) = 0$ & $\Re(s) = 3/2$ (um 1 verschoben) \\
$1 - 4^{1-s} = 0$ & $s = 1 + 2\pi i k/\ln 4$ (Gitterfaktor) \\
\bottomrule
\end{tabular}
\end{center}

Die Riemann-Hypothese ist in dieser Lesart eine Aussage \"uber das
Spektrum des $\mathbb{Z}^4$-Torus -- also \"uber ein geometrisches
Objekt des Rahmens, nicht blo\ss{} \"uber eine Zahlenreihe.

\subsection{Der Casimir-Punkt liegt symmetrisch zur kritischen Linie}
\label{sec:casimir}

\textbf{[K]} Dok.~314 rechnet die $\zeta$-regularisierte Vakuumenergie
$E = \pi\,\zeta_M(-1/2)$ und findet die Ordnungsumkehr: D4 minimiert die
Epstein-Zeta f\"ur $s > 0$, bei $s = 0$ sind alle Gitter gleich, und bei
$s = -1/2$ liegt $\mathbb{Z}^4$ tiefer ($-0{,}9321$ gegen $-0{,}8686$).

Der Casimir-Punkt $s = -1/2$ hat unter der Funktionalgleichung
$\zeta(s) = \chi(s)\zeta(1-s)$ den Spiegelpunkt $s = 3/2$. Die
Symmetrieachse:
\[
  \frac{-1/2 + 3/2}{2} \;=\; \frac12
  \qquad\text{-- exakt die kritische Linie.}
\]

Das physikalisch besetzte Intervall $[-1/2,\,3/2]$ hat die kritische
Linie als Mittelpunkt: Konsistenz zwischen Physik und Arithmetik.

\textbf{[X]} Ein Beweis folgt daraus nicht -- die kritische Linie ist
ohnehin durch die Funktionalgleichung fixiert. Ebenso wenig kann $\xi$
sie verschieben: eine St\"orung der Ordnung $\xi \sim 10^{-4}$ (auch mit
RG-Verst\"arker $100\xi$) trifft eine Achse nicht, die durch eine
funktionale Identit\"at festliegt.

\section{Die entscheidende Obstruktion: Weyl gegen Riemann-von Mangoldt}
\label{sec:weyl}

Naheliegend w\"are, einen Operator zu suchen, dessen Eigenwerte die
Nullstellen sind (Hilbert-P\'olya-Programm), und daf\"ur
$\text{T}^4/\mathbb{Z}_3$ geeignet zu deformieren. \textbf{Dieser
gesamte Zweig ist ausgeschlossen -- auf Satzebene, nicht als
Anschauung.}

\textbf{[Q]} \emph{Weyl-Gesetz.} F\"ur \emph{jede} kompakte Riemannsche
Mannigfaltigkeit beliebiger Dimension $d$ gilt asymptotisch
$N(\lambda) \sim C_d\,\text{Vol}\,\lambda^{d/2}$ -- eine reine Potenz.
Kr\"ummung, Rand und Orbifold-Punkte \"andern nur niedrigere Ordnungen,
nicht den Leitterm.

\textbf{[Q]} \emph{Riemann-von Mangoldt (bewiesen).}
\[
  N(T) = \frac{T}{2\pi}\ln\frac{T}{2\pi} - \frac{T}{2\pi}
  + \frac78 + O(1/T).
\]

\textbf{[K]} Der Vergleich \"uber den lokalen Exponenten
$\mathrm d\ln N/\mathrm d\ln T$:

\begin{center}
\begin{tabular}{rrr}
\toprule
$T$ & $N(T)$ & lokaler Exponent $d/2$ \\
\midrule
$10^{3}$ & $6{,}48\cdot10^{2}$ & $1{,}2457$ \\
$10^{5}$ & $1{,}38\cdot10^{5}$ & $1{,}1153$ \\
$10^{8}$ & $2{,}48\cdot10^{8}$ & $1{,}0642$ \\
$10^{12}$ & $3{,}95\cdot10^{12}$ & $1{,}0403$ \\
$10^{20}$ & --- & $1{,}0233$ \\
\bottomrule
\end{tabular}
\end{center}

Der Exponent ist $1 + 1/\ln(T/2\pi)$ -- er f\"allt monoton gegen 1, wird
aber \emph{nie konstant}. Spanne \"uber 17 Dekaden: $0{,}22$. Der beste
Potenzfit im Fenster $[10^3, 10^6]$ weicht systematisch um bis zu
$11\,\%$ ab.

\begin{keyresult}[Ausschluss auf Satzebene]
\textbf{[X]} Es gibt keine kompakte Riemannsche Mannigfaltigkeit --
gleich welcher Dimension, Kr\"ummung, Topologie oder Orbifold-Struktur
--, deren Laplace-Spektrum die Riemann-Nullstellen sind. Grund:
$T\ln T$ ist keine Potenz $T^{d/2}$.
\end{keyresult}

\textbf{Folge.} Der Zweig \enquote{$\text{T}^4/\mathbb{Z}_3$ geeignet
deformieren} ist vollst\"andig erledigt, nicht nur im flachen Fall. Das
zus\"atzliche $\ln T$ ist die Signatur eines
\textbf{Skalenfreiheitsgrads}: einer Dilatationsrichtung mit
logarithmisch wachsendem Volumen. Das ist kein Laplace auf einer
Mannigfaltigkeit.

\textbf{Zur\"uckgenommene Zwischendiagnose.} In einem ersten Durchgang
war argumentiert worden, $\log(p)$-Geod\"aten erforderten negative
Kr\"ummung. Das ist falsch: Was Log-L\"angen erzeugt, ist eine
Dilatationsstruktur (multiplikative Gruppe $\mathbb{R}_+^*$), nicht
Kr\"ummung; der Berry-Keating-Operator $H = xp$ ist der Erzeuger der
Dilatation, kein hyperbolischer Fluss. Die Weyl-Obstruktion ist die
korrekte und deutlich st\"arkere Aussage und macht die Kr\"ummungsfrage
gegenstandslos.

\section{Der konstruktive Anschluss -- und warum er nicht tr\"agt}
\label{sec:zelle}

Wenn die Nullstellendichte einen Skalenfreiheitsgrad verlangt, ist die
zust\"andige Dom\"ane die \textbf{ausgerollte} (Dok.~315) --
multiplikativ, logarithmisch, mit Dilatation als nat\"urlicher
Operation.

\subsection{Der Ansatz}

\textbf{[Q]} \emph{Berry-Keating (1999)}~\cite{berry1999}. Die semiklassische Abz\"ahlung
des Dilatationserzeugers $H = (xp+px)/2$ in einem Phasenraum mit
Minimalzelle:
\[
  N(E) = \frac{1}{2\pi}\bigl[\,E\ln\bigl(E/(l_x l_p)\bigr)
  - E + l_x l_p\,\bigr].
\]
Mit $l_x l_p = 2\pi$ wird daraus exakt die
Riemann-von-Mangoldt-Formel (verifiziert: Differenz konstant $0{,}125 =
1 - 7/8$, reiner Randterm, \"uber $T = 15 \ldots 5000$).

Naheliegend w\"are nun die Kette: Dok.~314 Kap.~D1 f\"uhrt
$\lambda_4\cdot m = 2\pi$ als exakte Ort-Relation; in der ausgerollten
Dom\"ane ist die zur Skala konjugierte Gr\"o\ss e der Dilatationsimpuls;
also stelle die FFGFT die Zellbedingung bereit.

\subsection{Warum die Kette zerf\"allt}

\textbf{[X] Stufe 1 -- die Zellbedingung ist inhaltsleer.}
$l_x l_p = 2\pi\hbar$ ist in nat\"urlichen Einheiten die Planck-Zelle
und folgt aus der kanonischen Quantisierung \emph{jedes} konjugierten
Paares. Dass $\lambda_4 m = 2\pi$ dieselbe Form hat, ist kein
eigenst\"andiger Befund -- die de-Broglie-Relation \emph{ist} die
Aussage $\hbar = 1$. Hinzu kommt: die Z\"ahlung ist gegen\"uber der
Zelle logarithmisch unempfindlich (Faktor 2 falsch $\Rightarrow$ $N$
\"andert sich um ${\sim}1\,\%$); ein \enquote{Treffer} w\"are ohnehin
kein Test.

\textbf{[X] Stufe 2 -- der eigentliche Gehalt scheitert quantitativ.}
Berry-Keating verlangt sch\"arfer als die Zelle: zwei \emph{unabh\"angige}
Abschneidungen $x > l_x$ und $p > l_p$, deren Produkt die Planck-Zelle
ist. Der FFGFT-Doppel-Cutoff macht das pr\"ufbar:

\begin{center}
\begin{tabular}{lr}
\toprule
Gr\"o\ss e & Wert \\
\midrule
UV: $L_0 = \xi\,\ell_P$ & $2{,}155\cdot10^{-39}$ m \\
IR: $L^* = 4\bar\lambda_e/\xi^{10}$ & $8{,}698\cdot10^{26}$ m \\
Verh\"altnis $L^*/L_0$ & $4{,}04\cdot10^{65}$ \\
verlangt w\"are & $2\pi \approx 6{,}28$ \\
\midrule
Fehlbetrag & \textbf{65 Dekaden} \\
\bottomrule
\end{tabular}
\end{center}

\textbf{[X] Stufe 3 -- endliches Fenster gegen unendlich viele
Nullstellen.} $\ln(L^*/L_0) = 151{,}06$, das sind $16{,}93$
$\xi$-Leiterstufen; die Berry-Keating-Abz\"ahlung im Fenster gibt
${\sim}1816$ Moden. Ein endliches Skalenfenster tr\"agt endlich viele
Moden, die Nullstellen sind abz\"ahlbar unendlich. Selbst wenn Stufe~1
und~2 hielten, k\"ame h\"ochstens ein Anfangsabschnitt heraus -- und ein
Anfangsabschnitt ist keine Spektralidentit\"at.

\emph{Randbeobachtung, ausdr\"ucklich nicht gebucht:} $16{,}93$ liegt
nahe 17. Beide Cutoffs tragen eigene Unsicherheiten -- der Vorfaktor~4
ist per R73 offen, $\ell_P$ h\"angt an $G$. Ohne Fehlerbudget nach P35
nicht auswertbar.

\subsection{Was bleibt}

\textbf{Zur\"uckzunehmen} ist die Formulierung \enquote{die
Nullstellendichte folgt aus der FFGFT-Zellbedingung}. Sie suggeriert
einen Befund, wo eine Trivialit\"at steht.

\textbf{[B] Es bleibt eine Strukturaussage, und die ist nicht wenig:}
Dass die Z\"ahlfunktion \"uberhaupt die Form
$(T/2\pi)\ln(T/2\pi)$ hat, verlangt einen \textbf{Dilatationserzeuger}.
Ein Laplace auf einer kompakten Mannigfaltigkeit kann sie nicht liefern
(Abschnitt~\ref{sec:weyl}). Die inhaltliche Aussage ist die
\emph{Identifikation der Struktur} -- die ausgerollte Dom\"ane tr\"agt
die Dilatation --, nicht der Zahlenwert $2\pi$.

\section{Die fehlende Randbedingung -- Test negativ}
\label{sec:rand}

$H = xp$ hat \emph{kontinuierliches} Spektrum. Die Z\"ahlfunktion legt
die mittlere Dichte fest, nicht die einzelnen Nullstellen. Es fehlt die
Randbedingung, die aus dem Kontinuum ein diskretes Spektrum macht.

\textbf{Der nat\"urliche Kandidat.} Die $\xi$-Leiter erzeugt eine
diskrete Untergruppe $\Gamma_\xi = \{\xi^n\} < \mathbb{R}_+^*$; der
Quotient $\mathbb{R}_+^*/\Gamma_\xi$ ist ein \emph{Kreis} vom Umfang
$L = \ln(1/\xi) = 8{,}9227$. Zul\"assige Dilatationsfrequenzen w\"aren
Vielfache von $2\pi/L = 0{,}7042$.

\subsection{Testdisziplin}

\textbf{[K]} Bei N\"ahe-Argumenten sind drei Kontrollen n\"otig (P35):
Nullverteilung \"uber zuf\"allige Perioden im \emph{gleichen} Bereich;
Artefaktkontrolle, weil Perioden gr\"o\ss er als der Datenbereich
automatisch kleine Scores erzeugen; Out-of-Sample-Pr\"ufung mit
Look-elsewhere-Korrektur.

Die Artefaktkontrolle ist nicht theoretisch: Ein erster Testlauf meldete
ein scheinbares Signal bei Periode $2\pi/|\ln(74/75)| = 468$. Da alle
Nullstellen $< 102$ liegen, runden alle Quotienten auf 0 -- reines
Artefakt einer Nullverteilung, die den Bereich nicht abdeckte.
Dokumentiert, weil genau diese Fehlerart bei $\xi$-N\"ahe-Argumenten
wiederkehrt.

\subsection{Ergebnis}

\textbf{[K]} Kandidaten im zul\"assigen Bereich (Nullverteilung: Mittel
$0{,}2502$):

\begin{center}
\begin{tabular}{lrrr}
\toprule
Kandidat & Periode & Score & $p$-Wert \\
\midrule
$2\pi/\ln(1/\xi)$ & $0{,}7042$ & $0{,}2812$ & $0{,}92$ \\
$2\pi/\ln 75$ & $1{,}4553$ & $0{,}2043$ & $\mathbf{0{,}005}$ \\
$\ln(1/\xi)$ & $8{,}9227$ & $0{,}2891$ & $0{,}94$ \\
$2\pi$ & $6{,}2832$ & $0{,}2366$ & $0{,}18$ \\
\bottomrule
\end{tabular}
\end{center}

Der Grenzkandidat $2\pi/\ln 75$ wurde out-of-sample gepr\"uft
(Nullstellen 31--50): Score $0{,}2686$, $p = 0{,}81$;
look-elsewhere-korrigiert $p = 1{,}0$.

\begin{keyresult}[Randbedingung]
\textbf{[X]} Die $\xi$-Leiter liefert die fehlende Randbedingung nicht.
Kein $\xi$-abgeleiteter Periodenkandidat \"uberlebt die
Out-of-Sample-Pr\"ufung.
\end{keyresult}

\subsection{Warum das strukturell erwartbar war}

\textbf{[K]} Eine Ein-Perioden-Kompaktifizierung erzeugt ein
\emph{\"aquidistantes} Gitter. Die Nullstellen zeigen aber
GUE-Statistik mit Niveauabsto\ss ung (Montgomery-Odlyzko~\cite{montgomery1973,odlyzko1987}): mittlerer
normierter Abstand $0{,}9815$ (soll ${\sim}1$), Anteil Abst\"ande
$< 0{,}3$ gleich \textbf{$0{,}000$} -- gegen ${\sim}0{,}04$ bei GUE und
${\sim}0{,}26$ bei gitterartigem Spektrum. Ein Kreis kann GUE
prinzipiell nicht liefern.

Die fehlende Randbedingung tr\"agt die Arithmetik (Primzahlen), nicht
die Geometrie.

\section{Der harmonische Zugang: $\xi$ als Tonnetz-Punkt}
\label{sec:tonnetz}

Die bisherigen Abschnitte behandeln $\xi$ als \emph{Gr\"o\ss e} -- als
Parameter, dessen Kleinheit \"uber die Wirkung entscheidet. Ein zweiter
Zugang fragt stattdessen nach seiner \emph{harmonischen Struktur}, also
nach der Stellung von $\xi$ im Primzahlgitter. Er f\"uhrt an eine andere
Stelle der Theorie und liefert einen sch\"arferen Ausschluss.

\subsection{Die Ausgangsbeobachtung}

\textbf{[K]} Die Primfaktorzerlegung
\[
  \frac{1}{\xi} \;=\; 7500 \;=\; 2^2\cdot 3\cdot 5^4
\]
stellt $\xi$ als Punkt im Eulerschen Tonnetz dar, mit den Koordinaten
$(a_2, a_3, a_5) = (-2, -1, -4)$. Das ist dieselbe Gitterstruktur wie im
Kontrollfall der Musikspirale (Abschnitt~\ref{sec:faktor} folgend):
Primzahlen 2, 3, 5 als Achsen des 5-Limit-Raums. \"Aquivalent in
Log-Schreibweise:
\[
  \ln(1/\xi) \;=\; 2\ln 2 + \ln 3 + 4\ln 5 \;=\; 8{,}922658\ldots
\]

Die harmonische Lesart legt eine andere Anschlussstelle nahe als der
Laplace-Operator. Ein Tonnetz tr\"agt drei inkommensurable
Log-Frequenzen $\ln 2$, $\ln 3$, $\ln 5$; drei unabh\"angige Frequenzen
erzeugen ein quasiperiodisches, gerade nicht \"aquidistantes Spektrum.
Damit ber\"uhrt der Zugang die Arithmetik der $\zeta$-Funktion direkt.

\subsection{Die zust\"andige Anschlussstelle: die Weilsche Formel}

\textbf{[Q]} Die explizite Formel von Weil koppelt Nullstellen und
Primzahlen:
\[
  \sum_\rho h(\gamma) \;=\; (\text{archimedischer Term})
  \;-\; 2\sum_{p,k} \frac{\ln p}{p^{k/2}}\;\hat h(k\ln p).
\]
Das duale \emph{L\"angenspektrum} der Nullstellen ist damit
\[
  \mathcal{L} \;=\; \{\,k\ln p \;:\; p \text{ prim},\; k = 1,2,3,\ldots\,\},
\]
also Vielfache von Logarithmen \textbf{einzelner} Primzahlen. Ob $\xi$
dort auftaucht, ist unmittelbar pr\"ufbar: die Fouriertransformierte der
Nullstellendichte muss bei den L\"angen aus $\mathcal{L}$ Peaks zeigen.

\subsection{Messung}

\textbf{[K]} Getestet wird
$F(u) = \bigl|\sum_n e^{i\gamma_n u}\bigr|/N$ f\"ur die ersten 100
Nullstellen; als Hintergrund dienen 3000 zuf\"allige $u$ im selben
Bereich (Median $0{,}0304$, 95\,\%-Schwelle $0{,}2005$,
99\,\%-Schwelle $0{,}2572$).

\begin{center}
\begin{tabular}{rrl}
\toprule
$u$ & $F(u)$ & Zuordnung \\
\midrule
$1{,}09861$ & $0{,}2321$ & $\ln 3$ \quad(Peak, 95\,\%) \\
$1{,}60944$ & $0{,}2651$ & $\ln 5$ \quad(Peak, 99\,\%) \\
$1{,}94591$ & $0{,}2625$ & $\ln 7$ \quad(Peak, 99\,\%) \\
$2{,}39790$ & $0{,}2610$ & $\ln 11$ \quad(Peak, 99\,\%) \\
$2{,}56495$ & $0{,}2676$ & $\ln 13$ \quad(Peak, 99\,\%) \\
\bottomrule
\end{tabular}
\end{center}

Die Primzahlen sind sichtbar; die Methode arbeitet. Dieselbe Messung
f\"ur $\xi$-abgeleitete L\"angen:

\begin{center}
\begin{tabular}{rrl}
\toprule
$u$ & $F(u)$ & Gr\"o\ss e \\
\midrule
$8{,}92266$ & $0{,}1775$ & $\ln(1/\xi) = 2\ln2+\ln3+4\ln5$ \\
$4{,}46133$ & $0{,}0777$ & $\ln(1/\xi)/2$ \\
$4{,}31749$ & $0{,}0384$ & $\ln 75 = \ln3+2\ln5$ \\
$2{,}23066$ & $0{,}0189$ & $\ln(7500)/4$ \\
$0{,}70418$ & $0{,}1313$ & $2\pi/\ln(1/\xi)$ \\
\bottomrule
\end{tabular}
\end{center}

Keine erreicht die 95\,\%-Schwelle.

\textbf{Aufl\"osungsgrenze.} Die Messung ist erst auswertbar, wenn $u$
die Phasen \"uber den vollen Kreis streut, also
$u \gg 2\pi/(\gamma_{\max}-\gamma_{\min}) = 0{,}02824$. F\"ur $u\to 0$
laufen alle Phasen zusammen und $F(u)\to 1$, unabh\"angig von jeder
Struktur. Kandidaten unterhalb dieser Grenze -- etwa
$\ln(75/74) = 0{,}01342$ -- sind mit 100 Nullstellen \emph{nicht
messbar}; sie liefern weder Best\"atigung noch Widerlegung.

\subsection{Der strukturelle Grund}

\begin{keyresult}[Kombination gegen Faktorisierung]
Tonnetz und Euler-Produkt teilen die Primzahlachsen, nutzen sie aber
\textbf{entgegengesetzt}. Das Euler-Produkt
$\zeta(s) = \prod_p (1-p^{-s})^{-1}$ logarithmiert zu einer Summe
\"uber Primzahlpotenzen: jede Primzahl steht f\"ur sich, ohne
Mischterme -- daher enth\"alt $\mathcal{L}$ nur $k\ln p$. Das Tonnetz
lebt vom Gegenteil: ein Intervall ist $2^a 3^b 5^c$ mit mehreren
nichtnullen Exponenten, und gerade die Mischung erzeugt die Kommas.
$\xi$ ist ein solcher Mischpunkt.
\end{keyresult}

\textbf{[X]} $\ln(1/\xi)$ ist damit eine \emph{Summe} von L\"angen und
keine L\"ange; Summen erscheinen in der Spurformel nicht als Peaks. Die
n\"achste echte L\"ange ist $3\ln 19 = 8{,}83332$ im Abstand
$0{,}08934$ -- ohne Bedeutung, da die Nachbarschaft im dichten
L\"angenspektrum nichts aussagt.

\textbf{Sch\"arfe des Ausschlusses.} Er ist strukturell, nicht
gr\"o\ss enordnungsbedingt: Er gilt unver\"andert, wenn $\xi$ von
Ordnung 1 w\"are. Was z\"ahlt, ist nicht der Wert von $\xi$, sondern
dass $1/\xi$ zusammengesetzt ist. Damit tr\"agt dieser Abschnitt den
st\"arkeren Grund als jede Absch\"atzung \"uber die Kleinheit von $\xi$.

\textbf{Verbleibende Pr\"uffl\"ache.} Die \emph{einzelnen} Achsen
$\ln 2$, $\ln 3$, $\ln 5$ sind L\"angen im Weil-Spektrum -- aber sie
sind es f\"ur jede Theorie, die \"uberhaupt Primzahlen enth\"alt; die
Zerlegung von 7500 selegiert nichts. Eine Aussage entst\"unde erst,
wenn die \emph{Exponenten} $(2,1,4)$ eine Gewichtung im Spektrum
erzeugten. Der folgende Abschnitt pr\"uft das.

\subsection{Schlie\ss t das 5-Limit die h\"oheren ein?}
\label{sec:einschluss}

Die Limit-Erweiterung von Test~B wirft die Umkehrfrage auf: L\"asst sich
zeigen, dass das 5-Limit die h\"oheren Limits bereits enth\"alt? Die
Antwort ist f\"unffach gestaffelt, und die Staffelung wiederholt das
Muster des gesamten Dokuments
(\texttt{z9\_limit\_einschluss.py}).

\textbf{[K] (1) Exakt: nein, beweisbar.} $\ln 7 = a\ln2 + b\ln3 +
c\ln5$ mit rationalen $a,b,c$ hie\ss e durchmultipliziert
$7^q = 2^A 3^B 5^C$ -- Widerspruch zur Eindeutigkeit der
Primfaktorzerlegung. Das gilt f\"ur jede Primzahl oberhalb des Limits,
also f\"ur unendlich viele.

\textbf{[K] (2) Als dichte N\"aherung: ja, beliebig genau.} Da
$\ln2, \ln3, \ln5$ \"uber $\mathbb{Q}$ linear unabh\"angig sind, liegt
das Gitter dicht in $\mathbb{R}$:

\begin{center}
\begin{tabular}{rrrlr}
\toprule
$N$ & Fehler & Cent & Exponenten & Tenney \\
\midrule
$2$ & $2{,}82\cdot10^{-2}$ & $48{,}77$ & $(2,2,-1)$ & $5{,}6\cdot10^{-3}$ \\
$4$ & $7{,}97\cdot10^{-3}$ & $13{,}80$ & $(-1,-2,3)$ & $4{,}4\cdot10^{-4}$ \\
$6$ & $4{,}45\cdot10^{-3}$ & $7{,}71$ & $(-5,2,2)$ & $1{,}4\cdot10^{-4}$ \\
$8$ & $2{,}29\cdot10^{-4}$ & $0{,}40$ & $(1,7,-4)$ & $3{,}7\cdot10^{-7}$ \\
\bottomrule
\end{tabular}
\end{center}

Der Fehler der besten N\"aherung $(1,7,-4)$ betr\"agt exakt
$\ln(4375/4374)$ -- \textbf{das Ragisma}. Die grobe N\"aherung
$(6,-2,0) = \ln(64/9)$ hat als Fehler $\ln(64/63)$, \textbf{das
Archytas-Komma}. Beide sind aus dem Euler-Kontrollfall bekannt: dort
als Beinaheschl\"usse im 7-Limit, hier als Fehler der
5-Limit-N\"aherung an die 7-Achse. Es ist dieselbe Gr\"o\ss e aus der
Gegenrichtung gelesen.

\textbf{[K] (3) Bei endlicher Aufl\"osung: ja, ununterscheidbar.} Die
Peakbreite der Messung aus Abschnitt~\ref{sec:bikoh} betr\"agt
$2\pi/(\gamma_{\max}-\gamma_{\min}) = 0{,}0021$; das Ragisma ist mit
$2{,}29\cdot10^{-4}$ nur $0{,}11$ Peakbreiten gro\ss{}:

\begin{center}
\begin{tabular}{lrr}
\toprule
Ort & $u$ & Leistung \\
\midrule
$\ln 7$ exakt & $1{,}945910$ & $2{,}144$ \\
N\"aherung $(1,7,-4)$ & $1{,}945682$ & $2{,}144$ \\
grob $(6,-2,0) = \ln(64/9)$ & $1{,}961659$ & $1{,}417$ \\
Kontrolle $\ln 7 + 0{,}05$ & $1{,}995910$ & $0{,}042$ \\
\bottomrule
\end{tabular}
\end{center}

Die feine N\"aherung ist von $\ln 7$ nicht zu unterscheiden, die grobe
f\"allt ab, der Kontrollpunkt liegt im Rauschen. Bei dieser Aufl\"osung
enth\"alt das 5-Limit die 7-Achse also effektiv -- eine Aussage \"uber
das Messverfahren, nicht \"uber die Arithmetik.

\textbf{[K] (4) Amplitudengewichtet: nein, und deutlich.} Das
Tenney-Gewicht der N\"aherung $(1,7,-4)$ betr\"agt
$3{,}7\cdot10^{-7}$ gegen $1/7$ f\"ur $\ln 7$ als eigene Achse --
Faktor $3{,}9\cdot10^{5}$. Das erkl\"art zugleich, \emph{warum} die
N\"aherung so genau sein kann: Sie braucht gro\ss e Exponenten, und
gro\ss e Exponenten bedeuten kleines Gewicht. Genauigkeit und Gewicht
sind gegenl\"aufig.

\textbf{[B] (5) Durch Temperierung: ja, aber per Setzung.} Erkl\"art
man das Ragisma zu $1$, so \emph{ist} $7 = 2\cdot 3^7/5^4$ im
temperierten System; ebenso macht das Wegtemperieren des
Archytas-Kommas die N\"aherung $64/9$ exakt.

\begin{keyresult}[Einschlie\ss ung entsteht nicht aus Struktur]
Das 5-Limit schlie\ss t die h\"oheren Limits nicht ein -- au\ss er man
l\"asst N\"aherung, begrenzte Aufl\"osung oder Temperierung als
Einschlie\ss ung gelten. Das ist dasselbe Muster wie \"uberall in
diesem Dokument: Einschlie\ss ung entsteht durch \emph{Rationalisierung}
(Temperierung) oder durch \emph{Unsch\"arfe} (Aufl\"osung), nie durch
Struktur. Ebenso schlie\ss t die Euler-Spirale nur durch Temperierung,
und ebenso sitzt Kopplung im Weil-Spektrum nur bei exakten
Primzahlpotenzen (Abschnitt~\ref{sec:bikoh}).
\end{keyresult}

F\"ur $\xi$ folgt nichts Neues, aber eine Best\"atigung: Dass $\xi$ ein
reiner 5-Limit-Punkt ist, l\"asst sich durch keine Erweiterung heilen --
und die N\"aherungen, die eine Erweiterung ersetzen k\"onnten, tragen
kein Gewicht.

\subsection{Messung statt Behauptung: der Bikoh\"arenztest}
\label{sec:bikoh}

Die Aussage, das Euler-Produkt besitze keine Mischterme, war bis hierher
aus der \emph{Form} des Produkts begr\"undet. Sie l\"asst sich messen --
mit dem daf\"ur zust\"andigen Verfahren und mit eingebauten Kontrollen.

\textbf{Aufbau.} Die Nullstellenfolge wird als Punktma\ss{} aufgefasst;
ihre Transformierte in der L\"angenvariablen ist
$F(u) = \sum_n e^{i\gamma_n u}$. \"Uber H\"ohenbl\"ocke gemittelt
ergeben sich Bispektrum und Bikoh\"arenz nach
Abschnitt~\ref{sec:quellen}. Die Nullstellen werden per
Riemann-Siegel-Formel selbst berechnet: 2469 St\"uck bis $t = 3000$
(Riemann-von Mangoldt erwartet $2468{,}6$), Restfehler
$7{,}5\cdot10^{-3}$ in $\gamma$, entsprechend $0{,}02$~rad
Phasenfehler bei $u \approx 2{,}6$.

\textbf{Testplan.} Das Weil-Spektrum enth\"alt $k\ln p$. Daraus folgen
Vorhersagen mit Positiv- \emph{und} Negativkontrollen: Ist die Summe
zweier L\"angen wieder eine L\"ange -- also eine Primzahlpotenz --, muss
Kopplung auftreten; f\"uhrt sie auf eine zusammengesetzte Zahl, darf
keine auftreten.

\textbf{Vorkontrolle.} Die Leistung von $F(u)$ zeigt die Einzell\"angen
$\ln 2 \ldots \ln 13$ s\"amtlich \"uber der 99\,\%-Schwelle. Der
Hintergrund muss dabei um die 13 echten L\"angen im Fenster bereinigt
werden -- zuf\"allige $u$ treffen mit etwa 2\,\% Wahrscheinlichkeit eine
davon und bilden sonst den Ausrei\ss erschwanz (Median $0{,}055$ gegen
$99\,\%$-Schwelle $0{,}881$ nach Bereinigung).

\textbf{Ergebnis} (24 H\"ohenbl\"ocke zu je 102 Nullstellen; kritischer
Wert aus 2000 phasenrandomisierten Realisierungen je Paar):

\begin{center}
\begin{tabular}{llrrll}
\toprule
Paar & Summe & $b^2$ & $b_c(99\%)$ & Urteil & erwartet \\
\midrule
$(\ln2,\ln2)$ & $\ln 4 = 2^2$ & $\mathbf{0{,}8088}$ & $0{,}2612$ & Kopplung & Kopplung \\
$(\ln3,\ln3)$ & $\ln 9 = 3^2$ & $\mathbf{0{,}8462}$ & $0{,}2214$ & Kopplung & Kopplung \\
$(\ln2,2\ln2)$ & $\ln 8 = 2^3$ & $\mathbf{0{,}7586}$ & $0{,}2432$ & Kopplung & Kopplung \\
$(\ln2,\ln3)$ & $\ln 6$ zus. & $0{,}0253$ & $0{,}1832$ & keine & keine \\
$(\ln2,\ln5)$ & $\ln 10$ zus. & $0{,}0252$ & $0{,}1912$ & keine & keine \\
$(\ln3,\ln5)$ & $\ln 15$ zus. & $0{,}0175$ & $0{,}1718$ & keine & keine \\
\bottomrule
\end{tabular}
\end{center}

\begin{keyresult}[Mischterme: gemessen]
Sechs von sechs Vorhersagen getroffen, mit einem Kontrast von Faktor
30 bis 50 zwischen Primzahlpotenzen und zusammengesetzten Zahlen.
Kopplung tritt genau dort auf, wo die Summe zweier L\"angen wieder eine
L\"ange ist, und nirgends sonst. \textbf{[K]} Die Mischterm-Aussage ist
damit gemessen, nicht nur strukturell begr\"undet.
\end{keyresult}

\textbf{Folge f\"ur $\xi$.} $\ln(1/\xi) = 2\ln2 + \ln3 + 4\ln5$ ist eine
Summe \"uber \emph{drei verschiedene} Primzahlen und f\"uhrt auf die
zusammengesetzte Zahl 7500. Nach dem hier gemessenen Muster kann dort
keine Kopplung sitzen -- was der unabh\"angige Fourier-Befund aus
Abschnitt~\ref{sec:tonnetz} ($F = 0{,}1775$, unter jeder Schwelle)
best\"atigt. Der offene Punkt Z-4 ist damit geschlossen, und der
Ausschluss aus Abschnitt~\ref{sec:tonnetz} ruht nun auf einer Messung
statt auf einem Formargument.

\subsection{Kann eine Gewichtung existieren? Drei Tests}
\label{sec:gewichtung}

Die verbleibende Pr\"uffl\"ache lautete: Eine Aussage entst\"unde erst,
wenn die Exponenten $(2,1,4)$ eine Gewichtung im Weil-Spektrum
erzeugten. Drei unabh\"angige Tests entscheiden diese Frage negativ
(\texttt{z7\_gewichtung\_unmoeglich.py}).

\textbf{[K] Test A -- Gewichtungsst\"arke.} H\"ohere Harmonische haben
kleinere Amplitude; das Standardma\ss{} daf\"ur ist die Tenney-H\"ohe,
Gewicht $\sim 1/(nd)$. Mit einem Exponenten wird die St\"arke frei:
\[
  w \;=\; \bigl(2^{|a|}\,3^{|b|}\,5^{|c|}\bigr)^{-\alpha},
\]
mit $\alpha = 0$ ungewichtet und $\alpha = 1$ Tenney. Die Gewichtssumme
hat die geschlossene Form
$\prod_p \bigl[1 + 2p^{-\alpha}/(1-p^{-\alpha})\bigr]$ und konvergiert
f\"ur jedes $\alpha > 0$:

\begin{center}
\begin{tabular}{rrrl}
\toprule
$\alpha$ & Gewichtssumme & eff.\ Punktzahl & Charakter \\
\midrule
$0{,}10$ & $4\,284$ & $23\,843$ & blind-dicht \\
$0{,}30$ & $245{,}8$ & $1\,767$ & Mitte \\
$0{,}50$ & $56{,}9$ & $360$ & Mitte \\
$1{,}00$ & $9{,}0$ & $36$ & achsen-d\"unn \\
\bottomrule
\end{tabular}
\end{center}

Beide Extreme sind ungeeignet: ungewichtet trifft das Gitter alles und
selegiert nichts, bei Tenney bleiben nur die Achsen \"ubrig.
Entscheidend ist die Mitte bei $\alpha \approx 0{,}3$ mit einigen
hundert bis tausend tragenden Punkten. Der Selektionstest \"uber den
gesamten Bereich:

\begin{center}
\begin{tabular}{rrrrr}
\toprule
$\alpha$ & $M(\text{Weil})$ & $M(\text{Zufall})$ & Verh\"altnis & $p$ \\
\midrule
$0{,}15$ & $2{,}1205$ & $2{,}0943$ & $1{,}013$ & $0{,}27$ \\
$0{,}30$ & $0{,}5344$ & $0{,}5210$ & $1{,}026$ & $0{,}31$ \\
$0{,}50$ & $0{,}1414$ & $0{,}1493$ & $0{,}947$ & $0{,}64$ \\
$1{,}00$ & $0{,}0114$ & $0{,}0226$ & $0{,}503$ & $0{,}93$ \\
\bottomrule
\end{tabular}
\end{center}

Kein $\alpha$ erzeugt Selektion; das Verh\"altnis bleibt bei oder unter
eins. Die Gewichtungsst\"arke ist nicht das Problem.

\textbf{[K] Test B -- Limit-Erweiterung.} Das 5-Limit ist die einfache
Fassung des Tonnetzes; erweiterte Tonnetze nehmen weitere
Primzahlachsen auf, als n\"achste die 7. Getestet bei $\alpha = 0{,}3$;
Primzahlen bis zum jeweiligen Limit sind Achsen und daher als trivial
ausgeschlossen.

\begin{center}
\begin{tabular}{rrrrr}
\toprule
Limit & Achsen & Punkte & Verh\"altnis & $p$ \\
\midrule
$5$ & $3$ & $2\,197$ & $1{,}049$ & $0{,}28$ \\
$7$ & $4$ & $28\,561$ & $1{,}015$ & $0{,}38$ \\
$11$ & $5$ & $371\,293$ & $1{,}048$ & $0{,}11$ \\
\bottomrule
\end{tabular}
\end{center}

Schwerer als die $p$-Werte wiegt der strukturelle Punkt: jede
Erweiterung verschiebt nur die Grenze, sie schlie\ss t keine L\"ucke.

\begin{center}
\begin{tabular}{lll}
\toprule
Limit & trivial abgedeckt & fehlt vollst\"andig \\
\midrule
$5$ & $2, 3, 5$ & alle $p \ge 7$ \\
$7$ & $2, 3, 5, 7$ & alle $p \ge 11$ \\
$13$ & $2 \ldots 13$ & alle $p \ge 17$ \\
\bottomrule
\end{tabular}
\end{center}

Ein Tonnetz mit Limit $L$ deckt die endlich vielen Primzahlen
$p \le L$ per Definition ab und die unendlich vielen $p > L$
\"uberhaupt nicht -- dazwischen gibt es nichts. Die Weilsche Formel
summiert aber \"uber s\"amtliche Primzahlen. Der Grenzfall macht es
endg\"ultig: ein Tonnetz mit unendlichem Limit w\"are keine Struktur
mehr; es enthielte jede Primzahl als Achse, also genau das
Euler-Produkt selbst -- ohne Zusatzgehalt und ohne jeden Bezug zu
$\xi$.

Hinzu kommt: $\xi = 2^2/(2^4\cdot 3\cdot 5^4)$ hat auf allen Achsen ab
7 den Exponenten null. Eine Erweiterung vergr\"o\ss ert das Gitter,
\"andert aber die Stellung von $\xi$ darin nicht; $\xi$ bleibt ein
reiner 5-Limit-Punkt.

\textbf{[K] Test C -- Starrheit.} Dies ist der eigentliche Grund, und
er gilt unabh\"angig von jeder Zahlenrechnung. Die Gewichte der
expliziten Formel sind keine freien Parameter, sondern Konsequenz des
Euler-Produkts:
\[
  \log\zeta(s) = \sum_{p,k} \frac1k\,p^{-ks},
\]
woraus $\ln p / p^{k/2}$ durch Ableiten und Verschieben nach $s = 1/2$
folgt. Es gibt keine Stelle, an der ein zus\"atzlicher Faktor eingef\"ugt
werden k\"onnte, ohne das Euler-Produkt selbst zu \"andern. Setzt man
probeweise $\prod_p (1 - g\,p^{-s})^{-1}$ mit $g \neq 1$ an, entsteht
eine andere Dirichlet-Reihe mit verschobener Konvergenzabszisse und
\emph{ohne} Funktionalgleichung.

\begin{keyresult}[Warum keine Gewichtung m\"oglich ist]
Die Funktionalgleichung $\zeta(s) = \chi(s)\zeta(1-s)$ ist genau das,
was die kritische Linie definiert. Eine Gewichtung, welche die Lage
der Nullstellen erkl\"aren soll, zerst\"ort damit das Objekt, dessen
Nullstellen sie erkl\"aren will. \textbf{[X]}
\end{keyresult}

Test~C tr\"agt allein; A und B zeigen zus\"atzlich, dass auch die
numerischen Wege nichts hergeben -- weder \"uber die Wahl der
Gewichtung noch \"uber erweiterte Tonnetze. Der offene Punkt Z-1 ist damit nicht
unbeantwortet, sondern \textbf{negativ entschieden}.

\section{Bilanz}

\begin{center}
\begin{tabular}{ll}
\toprule
Aussage & Status \\
\midrule
$\zeta$ ist Faktor der Torus-Spektral-Zeta & \textbf{[K]} verifiziert \\
Casimir-Intervall hat kritische Linie als Achse & \textbf{[K]} \\
Nullstellendichte verlangt Dilatationserzeuger & \textbf{[B]} Struktur \\
Zellbedingung als eigenst\"andiger Befund & \textbf{[X]} nur $\hbar = 1$ \\
Doppel-Cutoff als BK-Abschneidung & \textbf{[X]} 65 Dekaden \\
Kompakte Mannigfaltigkeit mit Nullstellen-Spektrum & \textbf{[X]} Weyl \\
$\xi$-Leiter diskretisiert die Nullstellen & \textbf{[X]} Test negativ \\
$\ln(1/\xi)$ im Weil-L\"angenspektrum & \textbf{[X]} Summe, keine L\"ange \\
Mischterme im Euler-Produkt & \textbf{[K]} bikoh\"arenzgemessen \\
5-Limit enth\"alt h\"ohere Limits & \textbf{[X]} nur genaehert/temperiert \\
$K_\text{frak}$ als temperierendes Instrument & \textbf{[X]} kein Komma im Rest \\
Gewichtung durch die Exponenten $(2,1,4)$ & \textbf{[X]} Z-1 entschieden \\
$\xi$ verschiebt die kritische Linie & \textbf{[X]} funktional fix \\
FFGFT l\"ost die Riemann-Hypothese & \textbf{[X]} \\
\bottomrule
\end{tabular}
\end{center}

\textbf{Der ehrlichste Satz.} Von den drei gepr\"uften Verbindungen
tragen zwei als Identit\"aten \textbf{[K]} -- sie sind aber Konsistenz,
nicht Aussage \"uber die Lage der Nullstellen. Die dritte, konstruktive,
ist beim Ausformulieren zerfallen: was wie eine Herleitung aussah, ist
$\hbar = 1$. \"Ubrig bleibt eine Strukturaussage: die Nullstellendichte
verlangt einen Dilatationserzeuger, und der lebt in der ausgerollten
Dom\"ane.

\textbf{Was ein tragf\"ahiger Ansatz leisten m\"usste:} eine
nicht-triviale Randbedingung auf der ausgerollten Skalenrichtung, die
(a) unendlich viele Moden tr\"agt -- also kein endliches
Cutoff-Fenster ist -- und (b) GUE-Statistik statt \"Aquidistanz erzeugt.
Beide Forderungen zeigen in dieselbe Richtung: die Struktur muss
arithmetisch sein, nicht geometrisch.

\section{Mittelung, Temperierung, Aufl\"osung: wo sitzt das Instrument?}
\label{sec:instrument}

Die Euler-Spirale schlie\ss t nur durch Temperierung
(Abschnitt~\ref{sec:faktor}), und das 5-Limit enth\"alt die h\"oheren
Limits nur durch N\"aherung, Aufl\"osung oder Temperierung
(Abschnitt~\ref{sec:einschluss}). In beiden F\"allen entsteht die
Schlie\ss ung nicht aus der Struktur, sondern aus etwas, das von au\ss
en hinzutritt. Es lohnt, dieses Hinzutretende genauer zu bestimmen --
denn es sitzt in drei Bereichen an drei verschiedenen Stellen.

\subsection{Drei Bereiche, drei Orte}

\textbf{Beim Geh\"or} sitzt die Mittelung im \emph{Empf\"anger}:
kritische B\"ander, endliche Frequenzaufl\"osung, Verschmelzung nahe
beieinanderliegender T\"one. Das Messinstrument wird Teil des
Ergebnisses. Die Temperierung ist hier keine willk\"urliche Setzung,
sondern folgt der Eigenschaft des H\"orenden.

\textbf{In der reinen Mathematik} gibt es keinen Empf\"anger. Die
Verh\"altnisse bestehen unabh\"angig davon, ob jemand sie unterscheidet.
Soll etwas schlie\ss en, was von sich aus nicht schlie\ss t, muss die
Temperierung als \emph{Setzung} eingef\"ugt werden -- als bewusster
Eingriff, dokumentiert im Komma, das dabei verschwindet.

\textbf{In der Natur} ist offen, ob es einen Empf\"anger gibt. Damit
stellt sich die Frage, ob die fraktale Korrektur diese Rolle spielt.

\subsection{Die pr\"ufbare Fassung}

\textbf{[B]} Eine Temperierung hinterl\"asst stets einen Rest, und
dieser Rest ist ein \emph{Komma}: ein Verh\"altnis glatter Zahlen, also
von Produkten kleiner Primzahlen. Der Korpus kennt einen Rest -- die
$\approx 7$~eV aus P-315-2, das $45\,000$-fache des Messbodens. Die
Frage lautet damit pr\"ufbar: Hat dieser Rest Kommastruktur?

\textbf{[K] Er hat sie nicht.} Die besten rationalen N\"aherungen an
$1 + 1{,}36\cdot10^{-5}$ sind $73530/73529$, $882365/882353$ und
$1250017/1250000$; keine ist glatt. Die Kontrollprobe zeigt die
Trennsch\"arfe des Tests: Alle sechs gepr\"uften Kommas --
pythagoreisch, syntonisch, Schisma, Breedsma, Ragisma, Landscape --
werden als 5- oder 7-Limit-glatt erkannt.

\emph{Zur Testf\"uhrung:} Ohne eine Bedingung an die Trefferg\"ute
meldet der Glattheitstest die triviale N\"aherung $1/1$ als
\enquote{5-Limit-glatt} -- bei 100\,\% Restfehler. Die Bedingung ist
notwendig, nicht kosmetisch; sie geh\"ort zur selben Fehlerklasse wie
die Artefakte in den Abschnitten~\ref{sec:rand}
und~\ref{sec:tonnetz}.

\subsection{Wo die fraktale Korrektur steht}

\begin{keyresult}[K$_\text{frak}$ ist nicht das temperierende Instrument]
W\"are $K_\text{frak}$ eine Temperierung, m\"usste sie einen
kommaartigen Rest hinterlassen. Der vorhandene Rest hat diese Struktur
nicht. \textbf{[X]} Das pr\"azisiert P-315-2: Die $\approx7$~eV sind
eher Polmassen-Effekt oder echter Korrekturterm als Abfall einer
Rationalisierung.
\end{keyresult}

\textbf{[B]} Hinzu kommt eine Pr\"azisierung, die leicht \"ubersehen
wird: In FFGFT sitzt die Setzung nicht bei $K_\text{frak}$, sondern
eine Ebene tiefer -- bei $\xi = 4/30000$ selbst. Das ist der
rationale Eingriff; $K_\text{frak} = 74/75$ ist seine \emph{Folge},
nicht der Akt. Insofern ist die fraktale Korrektur keine Temperierung,
sondern die Buchhaltung eines bereits geschlossenen Zyklus.

\textbf{[B]} Gleichwohl sitzt sie strukturell dort, wo ein Instrument
s\"a\ss e: Dok.~314 Kap.~D2 f\"uhrt $K_\text{frak}$ als
\emph{Eigenschaft der Durchquerung}, die \"uber die Br\"ucken\-konstanten
hereinkommt, nicht \"uber den lokalen Operator. Das ist wortw\"ortlich
\enquote{das Messinstrument wird Teil des Ergebnisses}: $K_\text{frak}$
liegt nicht im Feld, sondern im Durchlaufen. Sie \"ahnelt einem
Messeffekt im \emph{Ort}, nicht im \emph{Charakter} -- sie ist exakt,
nicht mittelnd.

\subsection{Der eigentliche Kandidat: Abschneiden statt Mitteln}

\textbf{[B]} Wenn die Natur ein Instrument hat, ist es eher der
Doppel-Cutoff. Er erzeugt eine Aufl\"osungsgrenze, und genau dort wird
das Verfahren Teil des Ergebnisses -- ohne dass irgendwo temperiert
w\"urde. Abschnitt~\ref{sec:einschluss} zeigt es an einem konkreten
Fall: Bei einer Peakbreite von $0{,}0021$ ist die 5-Limit-N\"aherung an
$\ln 7$ vom exakten Wert nicht mehr zu unterscheiden, obwohl das
Ragisma zwischen ihnen \emph{vorhanden} bleibt.

\begin{highlight}[Der Unterschied in einem Satz]
Eine Temperierung \emph{ver\"andert} die Werte und hinterl\"asst ein
Komma. Eine Aufl\"osungsgrenze ver\"andert nichts und hinterl\"asst
nichts -- sie macht Unterschiede nur unsichtbar. Beide lassen
Unterscheidungen verschwinden; nur die erste erzeugt Reste.
\end{highlight}

\textbf{Damit die Dreiteilung.} Das Geh\"or mittelt \emph{aktiv} und
erzeugt Kommas. Die Mathematik hat keinen Empf\"anger; dort muss
temperiert werden, durch Setzung. Die Natur, sofern der Rahmen recht
hat, schneidet \emph{passiv} ab: Unterscheidungen verschwinden ohne
Rationalisierung -- und die fraktale Korrektur ist daran nicht
beteiligt, sondern die Skalengrenze.

\textbf{[B]} \emph{Status.} Der Kommatest ist ein Ausschluss und
schlie\ss t eine bestimmte Deutung von $K_\text{frak}$ aus. Er zeigt
nicht, dass in der Natur \"uberhaupt kein Mittelungsmechanismus wirkt
-- nur, dass die fraktale Korrektur keiner ist. Ob der Doppel-Cutoff
als Instrument mehr ist als eine Analogie, bleibt offen; er h\"angt
seinerseits an Gr\"o\ss en, die der Korpus als offen f\"uhrt (Vorfaktor
4 in $L^*$, R73).

\section{Quellen und methodische Einordnung}
\label{sec:quellen}

Die in diesem Dokument verwendeten Verfahren sind Standard, ihre
Fallstricke ebenfalls dokumentiert. Der folgende Abschnitt ordnet die
verwertbare Literatur zu und benennt, wo die hier gerechneten Tests
hinter dem Stand der Technik zur\"uckbleiben.

\subsection{Tragende Quellen}

\textbf{Berry \& Keating (1999)}~\cite{berry1999}. Grundlage von
Abschnitt~\ref{sec:zelle}. Der Vergleich der Z\"ahlfunktionen legt
nahe, dass die $t_n$ Eigenwerte eines unbekannten hermiteschen
Operators sind, gewonnen durch Quantisierung eines klassischen Systems
mit Hamiltonfunktion $H_{cl} = xp$. Die Arbeit formuliert zugleich die
Anforderung an ein solches System: die \enquote{Riemann-Dynamik}
m\"usste chaotisch sein und periodische Bahnen besitzen, \emph{deren
Perioden Vielfache von Logarithmen der Primzahlen sind}. Das ist genau
das L\"angenspektrum aus Abschnitt~\ref{sec:tonnetz} -- und die
Bedingung, an der das Tonnetz scheitert, da es Summen statt Vielfacher
liefert.

\textbf{Sierra (2019)}~\cite{sierra2019} \textbf{sowie Sierra \& Townsend
(2008)}~\cite{sierra2008}. Zeigen den Stand des Hilbert-P\'olya-Programms: spektrale
Realisierungen \"uber $xp$ mit Randwechselwirkung, \"uber
Landau-Niveaus, zuletzt \"uber einen masselosen Dirac-Fermion im
Rindler-Raum mit Deltapotentialen auf den quadratfreien ganzen Zahlen.
Bemerkenswert f\"ur die hiesige Fragestellung: Alle diese Konstruktionen
kodieren die Arithmetik \emph{explizit} (quadratfreie Zahlen,
Modulformen, Randfunktionen aus der Riemann-Siegel-Formel). Keine
gewinnt sie aus einer Geometrie. Das st\"utzt den Befund von
Abschnitt~\ref{sec:weyl} von der anderen Seite: Wo die Nullstellen
spektral realisiert werden, geschieht es durch Einbau der
Zahlentheorie, nicht durch Wahl eines Raumes.

\textbf{P\'apics (2012)}~\cite{papics2012}. Die
methodisch n\"achstliegende Arbeit zu Abschnitt~\ref{sec:gewichtung}.
Untersucht wird, ob gefundene Frequenzen in harmonischer Relation
zueinander stehen -- dasselbe Problem in der Asteroseismologie.
Zentrales Ergebnis: Kombinationsfrequenzen treten in Zufallsdaten mit
hoher Wahrscheinlichkeit auf, und mit steigender Datenqualit\"at wird
die Unterscheidung zwischen echten und Scheinkombinationen
\emph{schwieriger}, nicht leichter. Der brauchbare Diskriminator ist
nicht die Existenz einzelner Kombinationen, sondern die \emph{Anzahl}
der Kombinationen niedrigster Ordnung im Vergleich zu Simulationen mit
gleichen Eigenschaften. Genau dieses Vorgehen -- Vergleich gegen eine
simulierte Nullverteilung statt Bewertung von Einzeltreffern -- liegt
den Permutationstests dieses Dokuments zugrunde.

\textbf{P\'ol\'oskei et al.\ (2018)}~\cite{poloskei2018}. Behandelt
das Verfahren, das f\"ur die Frage nach harmonischer Kopplung
eigentlich zust\"andig ist. Die Bikoh\"arenz
\[
  b^2(f_1,f_2) = \frac{|B(f_1,f_2)|^2}
  {A\{|X(f_1)X(f_2)|^2\}\,A\{|X(f_1{+}f_2)|^2\}}
\]
misst nicht Amplituden, sondern \emph{Phasenkopplung}: ob
$f_3 = f_1 + f_2$ mit fester Phasenrelation
$\varphi_3 = \varphi_1 + \varphi_2 + \text{const}$ auftritt.
Entscheidend f\"ur die hiesige Verwendung ist die Warnung vor
Falschpositiven: Bei nichtstation\"aren oder kurzen Signalen liefert die
klassische Bikoh\"arenz hohe Werte \emph{ohne} jede Kopplung. Das
vorgeschlagene Gegenmittel ist ein Monte-Carlo-Verfahren mit den
gemessenen Fourier-Amplituden bei \emph{randomisierten Phasen}, daraus
die Nullverteilung $\rho(b)$ und ein kritischer Wert $b_c(\alpha)$.
Ferner wird der Zielkonflikt benannt: Die Zahl der Mittelungen l\"asst
sich nur auf Kosten der Frequenzaufl\"osung erh\"ohen
($\Delta f = f_s N/M$).

\subsection{Weitere verwertbare Verfahren}

Die folgenden Methoden unterschreiten das Fourier-Aufl\"osungslimit,
indem sie ein Modell annehmen (Summe von Exponentialen) statt blind zu
transformieren. Sie sind hier nicht eingesetzt worden, w\"aren aber die
n\"achste Ausbaustufe:

\begin{itemize}
  \item \textbf{Unterraumverfahren:} MUSIC (Schmidt), ESPRIT
    (Roy/Kailath), Min-Norm. Zerlegen die Kovarianzmatrix in Signal-
    und Rauschunterraum.
  \item \textbf{Matrix-Pencil}~\cite{hua1990} f\"ur ged\"ampfte und unged\"ampfte Sinusoide.
  \item \textbf{Iterative Verfahren:} RELAX, NOMP; sowie
    Volterra-Reihen-Ans\"atze zur selbstkonsistenten Entfernung von
    Kombinationsfrequenzen~\cite{lares2020}.
  \item \textbf{Multitaper} (Thomson) mit F-Test auf Linienkomponenten;
    \textbf{Multichannel Singular Spectrum Analysis} zur Extraktion
    dynamischer Moden ohne Fourier-Basis.
\end{itemize}

Zur Gr\"o\ss enordnung des Gewinns: In einem publizierten Vergleich
dreier Frequenzen bei 440/445/450~Hz \"uber 100~ms erreicht ESPRIT
einen Frequenz-RMSE von $0{,}087$~Hz, spektrales MUSIC dagegen
$2{,}88$~Hz -- der Unterschied zwischen den Modellklassen ist also
gr\"o\ss er als der zur Fourier-Analyse.

\subsection{Wo die hiesigen Tests zur\"uckbleiben}

\textbf{[B]} Die Skripte \texttt{z6} und \texttt{z7} verwenden die
nackte Fouriersumme
$F(u) = |\sum_n e^{i\gamma_n u}|/N$ -- das einfachste verf\"ugbare
Werkzeug. Drei Einschr\"ankungen folgen daraus:

\begin{enumerate}
  \item \emph{Keine Kopplungsmessung.} $F(u)$ misst, ob bei der
    L\"ange $u$ \"uberhaupt Struktur sitzt, nicht ob zwei L\"angen
    \emph{gekoppelt} sind. Die strukturelle Behauptung, das
    Euler-Produkt besitze keine Mischterme, ist in diesem Dokument
    aus der Form des Produkts begr\"undet. Der Bikoh\"arenztest ist
    inzwischen nachgeholt (Abschnitt~\ref{sec:bikoh}) und best\"atigt
    sie; die \"ubrigen Tests dieses Dokuments bleiben davon unber\"uhrt
    auf dem einfacheren Verfahren.
  \item \emph{Statistik.} Die Abschnitte \ref{sec:tonnetz} und
    \ref{sec:gewichtung} rechnen mit 100 Nullstellen aus der
    Literatur; f\"ur Bispektren ist das zu wenig. Abschnitt
    \ref{sec:bikoh} berechnet daher 2469 Nullstellen selbst.
  \item \emph{Kein Zeitsignal.} Die Nullstellenfolge ist kein
    Messsignal; Stationarit\"at ist nicht definiert, und die
    Standardtests auf Nichtstationarit\"at greifen nicht unmittelbar.
\end{enumerate}

Diese Einschr\"ankungen ber\"uhren die Befunde \textbf{[X]} dieses
Dokuments nicht: Die Weyl-Obstruktion
(Abschnitt~\ref{sec:weyl}) und die Starrheit der Weil-Gewichte
(Abschnitt~\ref{sec:gewichtung}, Test~C) sind strukturelle Argumente
und von der Wahl des Spektralverfahrens unabh\"angig. Sie betreffen
allein die numerischen Tests, deren Aussagekraft dadurch nach unten
begrenzt ist -- was den negativen Gesamtbefund nicht schw\"acht,
sondern nur seinen numerischen Anteil relativiert.

\section{Schlussbetrachtung: warum sich nicht zur\"uckrechnen l\"asst}

Am Ende steht ein Bild, das die einzelnen Befunde zusammenfasst, und
eine Deutung, die \"uber sie hinausgeht. Beides ist zu trennen.

\subsection{Was die Rechnungen zeigen}

Die Primzahlen sind Achsen -- im Tonnetz wie im Euler-Produkt.
Dieselben Zahlen, dieselbe Gitterstruktur. Bei einfachen
Verh\"altnissen ist der R\"uckschluss auf die zugrundeliegende
Harmonie eindeutig: $\ln 2$, $\ln 3$, $\ln 5$ erscheinen als klare
Linien im Spektrum der Nullstellen, jede an ihrer Stelle, jede
zuzuordnen.

Mit der Zusammensetzung geht diese Eindeutigkeit verloren, und zwar
aus zwei Gr\"unden, die beide in den Tests sichtbar werden.

Der erste ist die \emph{Mehrdeutigkeit}. Ein einfaches Verh\"altnis hat
wenige m\"ogliche Zerlegungen. Ein zusammengesetztes hat viele -- und
je weiter man das Gitter aufspannt, desto dichter liegen die
M\"oglichkeiten, bis schlie\ss lich jede Zahl irgendwie hineinpasst.
Nicht die Seltenheit einer Kombination macht die Schwierigkeit,
sondern ihre Vieldeutigkeit: Wo alles erreichbar ist, sagt das
Erreichen nichts mehr.

Der zweite Grund wiegt schwerer. Das Euler-Produkt behandelt jede
Primzahl \emph{f\"ur sich}, ohne Mischterme; sein Spektrum kennt nur
$k\ln p$, Vielfache einzelner Primlogarithmen. Sobald ein Verh\"altnis
mehrere Primzahlen mischt, ist es keine einzelne Harmonie mehr,
sondern eine \"Uberlagerung -- und \"Uberlagerungen erscheinen im
Spektrum nicht als eigene Linie. $\xi = 2^2/(2^4\cdot3\cdot5^4)$ ist
eine solche \"Uberlagerung.

Ein Bild aus der Praxis trifft es genau: Bei einer einzelnen Zunge
l\"asst sich aus dem Klang auf den Grundton schlie\ss en. Bei einem
vollen Akkord \"uber mehrere Register h\"ort man die Mischung -- aber
der Grundton der Mischung existiert nicht. Es gibt nur die Grundt\"one
der einzelnen Zungen. $\xi$ ist ein Akkord, kein Ton.

\subsection{Die Richtung, die geht -- und die, die nicht geht}

Damit l\"asst sich pr\"azise sagen, was hier eigentlich verschlossen
ist. Es ist nicht die Verbindung zwischen Harmonie und Struktur; es
ist ihre \emph{Umkehrbarkeit}.

Vorw\"arts tr\"agt sie: Aus einfachen Verh\"altnissen l\"asst sich
Zusammengesetztes aufbauen, Schritt f\"ur Schritt, jeder Schritt
nachvollziehbar. So entstehen Intervalle aus Quinten und Terzen, so
entsteht $\xi$ aus $2$, $3$ und $5$.

R\"uckw\"arts tr\"agt sie nicht. Die Zusammensetzung ist nicht
umkehrbar, weil sie Information vernichtet: viele Wege f\"uhren an
dieselbe Stelle, und der zur\"uckgelegte Weg ist am Ergebnis nicht
mehr abzulesen. Das ist keine Schw\"ache der Methode, die sich mit
besseren Verfahren beheben lie\ss e, sondern eine Eigenschaft der
Sache. Die Tests dieses Dokuments haben es an drei verschiedenen
Stellen best\"atigt: bei der Gewichtungsst\"arke, bei der
Limit-Erweiterung und bei der Starrheit der Weil-Gewichte.

\subsection{Deutung}

\begin{highlight}[Deutung, kein Befund]
Es liegt nahe, dass komplexe Gebilde in der Natur auf solchen
Harmonien beruhen -- dass sich Zusammengesetztes aus einfachen
Verh\"altnissen aufbaut, wie ein Akkord aus T\"onen. Diese Vermutung
ist alt und in vieler Hinsicht fruchtbar gewesen; sie steht auch
hinter der Konstruktion von $\xi$.

Was dieses Dokument zeigt, ist etwas anderes und Bescheideneres: Selbst
wenn es so ist, l\"asst sich der Weg nicht zur\"uckgehen. Aus dem
fertigen Gebilde die zugrundeliegenden Harmonien zu rekonstruieren,
ist nicht schwer, sondern unm\"oglich -- die Zusammensetzung hat die
Zuordnung gel\"oscht.
\end{highlight}

Der erste Satz ist eine Deutung und bleibt es; er wird durch nichts in
diesem Dokument gest\"utzt und durch nichts widerlegt. Der zweite ist
gerechnet.

Zwischen beiden liegt der eigentliche Ertrag: Wer eine harmonische
Ordnung hinter komplexen Strukturen vermutet, muss sie \emph{vorw\"arts}
belegen -- durch Konstruktion und Vorhersage. Der R\"uckweg, aus der
fertigen Struktur auf die Ordnung zu schlie\ss en, steht nicht offen.
Das gilt f\"ur die Riemann-Hypothese, und es gilt allgemeiner.

\section{Offene Punkte}

\begin{itemize}
  \item \textbf{Z-1 (entschieden, negativ):} Eine Gewichtung des
    Weil-Spektrums durch die Tonnetz-Exponenten $(2,1,4)$ kann es
    nicht geben -- siehe Abschnitt~\ref{sec:gewichtung}. Das Tonnetz
    Weder eine Gewichtungsst\"arke noch eine
    Limit-Erweiterung erzeugt Selektion; vor allem aber sind die
    Gewichte der expliziten Formel keine freien Parameter.
  \item \textbf{Z-2:} Der Gitterfaktor $(1 - 4^{1-s})$ liefert
    Nullstellen bei $s = 1 + 2\pi i k/\ln 4$. Hat diese Leiter eine
    Korpus-Bedeutung (D4? Faktor 4 aus $L^* = 4\bar\lambda_e/\xi^{10}$,
    R73)? Bisher nicht gepr\"uft.
  \item \textbf{Z-3:} Der W\"armekern-Befund aus Dok.~314
    (D4/$\mathbb{Z}^4$-Unterschied exponentiell klein,
    $1{,}8\cdot10^{-13}$ bei $t = 0{,}1$) hat ein Gegenst\"uck in der
    $\zeta$-Theorie (Thetareihen-Feinstruktur). Ob daraus eine
    Pr\"uffl\"ache wird: offen.
  \item \textbf{Z-4 (geschlossen):} Der Bikoh\"arenztest ist
    durchgef\"uhrt (Abschnitt~\ref{sec:bikoh}) und best\"atigt die
    Mischterm-Aussage in sechs von sechs Kontrollen.
  \item \textbf{Nicht empfohlen:} weitere Suche nach $\xi$-Perioden in
    den Nullstellen. Abschnitt~\ref{sec:rand} schlie\ss t das mit dem
    GUE-Argument strukturell aus, nicht nur numerisch.
\end{itemize}

\subsection*{Anmerkung: bringen Quantenrechner hier eine Besserung?}
\addcontentsline{toc}{subsection}{Anmerkung: Quantenrechner}

\textbf{[B]} Die Frage liegt nahe, zumal Dok.~024 die
$\zeta$-Funktion im Zusammenhang mit Faktorisierung und dem
Shor-Algorithmus erw\"ahnt. Die Antwort ist negativ, und der Grund ist
derselbe wie in den Abschnitten zuvor.

Quantenrechner beschleunigen \emph{Suchprobleme}: Shor faktorisiert
schnell, Grover durchsucht schneller. Beides setzt voraus, dass es
etwas zu finden gibt -- einen Faktor, einen markierten Eintrag -- und
dass die Rechenzeit der Engpass ist. Genau das ist hier nicht der Fall:

\begin{itemize}
  \item Die Weil-Gewichte sind keine freien Parameter. Es gibt keinen
    Suchraum, in dem sich suchen lie\ss e; schnelleres Suchen nach
    etwas Nichtexistentem f\"uhrt zu nichts.
  \item Die Zusammensetzung vernichtet Information. Das ist kein
    Komplexit\"atsproblem, sondern Informationsverlust -- auch ein
    Quantenrechner rekonstruiert nicht, was nicht mehr vorhanden ist.
  \item Ein Tonnetz endlichen Limits l\"asst unendlich viele
    Primzahlen aus. Mehr Rechenleistung schiebt das Limit weiter und
    l\"asst immer noch unendlich viele aus.
\end{itemize}

Hinzu kommt die Grenze aller Numerik in dieser Frage: Nullstellen
lassen sich weiter hinaufrechnen als die derzeit bekannten
Gr\"o\ss enordnungen, aber jede weitere ist ein zus\"atzlicher
Datenpunkt f\"ur eine Aussage, die \"uber \emph{s\"amtliche}
Nullstellen l\"auft. Numerik kann die Riemann-Hypothese widerlegen --
ein Gegenbeispiel gen\"ugt --, beweisen kann sie sie nie.

Die Erw\"ahnung in Dok.~024 verbindet zwei Dinge, die assoziativ
zusammenh\"angen, aber sachlich nicht: das Faktorisieren einer
\emph{konkreten Zahl}, wo Quantenrechner tats\"achlich etwas leisten,
und eine Aussage \"uber die Struktur \emph{aller} Primzahlen, wo sie es
nicht tun.

\section{Verifikation}

Alle Zahlen sind durch zehn Skripte reproduzierbar (reine
Standardbibliothek, Sollwerte als Assertions, Seed 20780458):

\begin{itemize}
  \item \texttt{z1\_weyl\_obstruktion.py} -- Weyl-Gesetz gegen
    Riemann-von Mangoldt, lokaler Exponent, Potenzfit.
  \item \texttt{z2\_zellbedingung.py} -- Berry-Keating-Z\"ahlung,
    \"Aquivalenz zur Riemann-Form, Gegenprobe gegen Weyl.
  \item \texttt{z3\_randbedingung\_test.py} -- $\xi$-Leiter-Test mit
    Artefaktkontrolle, Nullverteilung, Out-of-Sample, GUE-Kontrolle.
  \item \texttt{z4\_epstein\_casimir.py} -- Faktorisierung gegen
    Direktsumme, Casimir-Symmetriepunkt.
  \item \texttt{z5\_zellbedingung\_kritik.py} -- Zerfall des
    Zellarguments in drei Stufen.
  \item \texttt{z6\_tonnetz\_weil.py} -- Tonnetz-Koordinaten von $\xi$,
    Fourier-Nachweis des Weil-L\"angenspektrums, Aufl\"osungsgrenze.
  \item \texttt{z10\_kommatest.py} -- hat der Restterm aus P-315-2
    Kommastruktur? Glattheitstest mit Kontrollprobe.
  \item \texttt{z9\_limit\_einschluss.py} -- f\"unf Lesarten der
    Limit-Einschlie\ss ung; Ragisma und Archytas-Komma als
    N\"aherungsfehler.
  \item \texttt{z8\_bikohaerenz.py} -- Nullstellen per Riemann-Siegel,
    Bikoh\"arenz mit Phasenrandomisierung, Positiv- und
    Negativkontrollen.
  \item \texttt{z7\_gewichtung\_unmoeglich.py} -- Gewichtungsst\"arke
    $\alpha$, Limit-Erweiterung (5-, 7-, 11-Limit), Permutationstests,
    Starrheit der Weil-Gewichte.
\end{itemize}

Eingaben: $\xi = 4/30000$ exakt; $\ell_P = 1{,}616255\cdot10^{-35}$~m,
$\bar\lambda_e = 3{,}8615926796\cdot10^{-13}$~m (CODATA~2022);
Nullstellenwerte (erste 50) aus der Literatur.

\begin{thebibliography}{99}

\bibitem{berry1999}
Berry, M.~V., \& Keating, J.~P. (1999).
\textit{The Riemann Zeros and Eigenvalue Asymptotics}.
SIAM Review, 41(2), 236--266.
\url{https://doi.org/10.1137/S0036144598347497}

\bibitem{sierra2019}
Sierra, G. (2019).
\textit{The Riemann Zeros as Spectrum and the Riemann Hypothesis}.
Symmetry, 11(4), 494.
\url{https://doi.org/10.3390/sym11040494}

\bibitem{sierra2008}
Sierra, G., \& Townsend, P.~K. (2008).
\textit{Landau Levels and Riemann Zeros}.
Physical Review Letters, 101, 110201.
\url{https://doi.org/10.1103/PhysRevLett.101.110201}

\bibitem{papics2012}
P\'apics, P.~I. (2012).
\textit{The puzzle of combination frequencies found in heat-driven
pulsators}.
Astronomische Nachrichten, 333(10), 1053--1056.
\url{https://doi.org/10.1002/asna.201211815}

\bibitem{poloskei2018}
P\'ol\'oskei, P.~Z., Papp, G., Por, G., Horvath, L., \& Pokol, G.~I.
(2018).
\textit{Bicoherence analysis of nonstationary and nonlinear processes}.
arXiv:1811.02973 [eess.SP].
\url{https://arxiv.org/abs/1811.02973}

\bibitem{kim1979}
Kim, Y.~C., \& Powers, E.~J. (1979).
\textit{Digital bispectral analysis and its applications to nonlinear
wave interactions}.
IEEE Transactions on Plasma Science, 7(2), 120--131.

\bibitem{hua1990}
Hua, Y., \& Sarkar, T.~K. (1990).
\textit{Matrix pencil method for estimating parameters of exponentially
damped/undamped sinusoids in noise}.
IEEE Transactions on Acoustics, Speech and Signal Processing, 38(5),
814--824.

\bibitem{lares2020}
Lares-Martiz, M., Garrido, R., \& Su\'arez, J.~C. (2020).
\textit{Self-consistent method to extract non-linearities from
pulsating star light curves -- I. Combination frequencies}.
Monthly Notices of the Royal Astronomical Society, 498(4), 6069--6086.
\url{https://doi.org/10.1093/mnras/staa2762}

\bibitem{euler1739}
Euler, L. (1739).
\textit{Tentamen novae theoriae musicae ex certissimis harmoniae
principiis dilucide expositae}.
Petropoli: Ex Typographia Academiae Scientiarum.

\bibitem{odlyzko1987}
Odlyzko, A.~M. (1987).
\textit{On the distribution of spacings between zeros of the zeta
function}.
Mathematics of Computation, 48(177), 273--308.

\bibitem{montgomery1973}
Montgomery, H.~L. (1973).
\textit{The pair correlation of zeros of the zeta function}.
Analytic Number Theory, Proc.\ Sympos.\ Pure Math.\ XXIV, 181--193.
American Mathematical Society.

\bibitem{codata2022}
Tiesinga, E., Mohr, P.~J., Newell, D.~B., \& Taylor, B.~N. (2024).
\textit{CODATA Recommended Values of the Fundamental Physical
Constants: 2022}.
Reviews of Modern Physics, 97, 025002.
\url{https://physics.nist.gov/cuu/Constants/}

\bibitem{pascher_ffgft}
Pascher, J. (2026).
\textit{FFGFT -- Fundamentale Fraktale Geometrische Feldtheorie}.
Korpus, insbesondere Dok.~230/231/232 (Hilbertraum), Dok.~295/313
(Schlie\ss ungsgabelung), Dok.~314 (D4-Gitter, Epstein-Zeta, Casimir),
Dok.~315 (eingerollt/ausgerollt), Dok.~190 (Korrekturregister).
\url{https://github.com/jpascher/T0-Time-Mass-Duality}

\end{thebibliography}
